\chapter{NFA Equivalence and Regular Expressions}

\section{Equivalence of DFA and NFA}

As established previously, Deterministic Finite State Automata (DFA) and Non-Deterministic Finite State Automata (NFA) differ in how their transition functions are defined. In a DFA, there is exactly one next state for each symbol read. In an NFA, a transition can lead to a \emph{set} of possible states, and acceptance is defined existentially (the word is accepted if \emph{at least one} computation path leads to a final state).

Despite these differences, DFA and NFA are perfectly equivalent in terms of expressive power. Any language that can be recognized by an NFA can also be recognized by a DFA. 

\subsection{The Power Set Construction}

To prove this equivalence constructively, we can convert any given NFA $\mathcal{A} = (Q, \mathcal{A}, \Delta, q_0, F)$ into an equivalent DFA $\mathcal{A}' = (Q', \mathcal{A}, \delta', q'_0, F')$ using the \textbf{power set construction}.

The components of the equivalent DFA are defined as follows:
\begin{enumerate}
\item \textbf{States ($Q'$)}: The states of the DFA are subsets of the original states $Q$. Thus, $Q' = 2^Q$. This highlights an exponential blow-up: a DFA may require exponentially more states than the corresponding NFA.
\item \textbf{Alphabet ($\mathcal{A}$)}: Remains unchanged.
\item \textbf{Initial State ($q'_0$)}: The singleton set containing the original initial state, $q'_0 = \{q_0\}$.
\item \textbf{Transition Function ($\delta'$)}: For any state $P \in Q'$ (where $P \subseteq Q$) and any symbol $a \in \mathcal{A}$, the transition is the union of all states reachable in the NFA from any state in $P$:
$$ \delta'(P, a) = \bigcup_{p \in P} \Delta(p, a) $$
\item \textbf{Final States ($F'$)}: A state $P \in Q'$ is considered final if it contains at least one state that was final in the original NFA:
$$ F' = \{ P \in Q' \mid P \cap F \neq \emptyset \} $$
\end{enumerate}

\section{NFA with $\varepsilon$-moves}

A useful extension to the NFA model is the \textbf{NFA with $\varepsilon$-moves}. This automaton functions like a standard NFA but has the additional ability to transition between states without consuming any input symbols (an $\varepsilon$-move). 

This is a powerful abstraction tool that allows us to ignore certain implementation details or "jump" between states freely. However, it is primarily useful for finite words; for infinite words, $\varepsilon$-moves lose their utility.

\subsection{$\varepsilon$-closure}

To handle these spontaneous transitions, we introduce the concept of the \textbf{$\varepsilon$-closure}. The $\varepsilon$-closure of a state $q$ is the set of all states that can be reached from $q$ by making zero or more $\varepsilon$-moves. 

For a set of states $P$, the $\varepsilon$-closure is simply the union of the $\varepsilon$-closures of each individual state $p \in P$.

\subsection{Extended Transition Function}

When processing a word $w$ and a symbol $a$, the extended transition function $\hat{\delta}(q, wa)$ for an NFA with $\varepsilon$-moves involves alternating between $\varepsilon$-closures and symbol consumption:
\begin{enumerate}
\item Compute the set of states reachable after reading $w$.
\item Take the $\varepsilon$-closure of those states.
\item Consume the symbol $a$ to reach a new set of states.
\item Take the $\varepsilon$-closure of this newly reached set of states.
\end{enumerate}

\subsection{Equivalence Theorem}

\begin{theorem}
For every NFA with $\varepsilon$-moves, there exists an equivalent NFA without $\varepsilon$-moves that recognizes exactly the same language.
\end{theorem}

To convert an NFA with $\varepsilon$-moves into a standard NFA, we define a new transition function where a single step consuming a symbol $a$ encapsulates the $\varepsilon$-closure before the move, the move itself, and the $\varepsilon$-closure after the move. The initial state is added to the final states if the $\varepsilon$-closure of the initial state contains a final state.

Ultimately, DFA, NFA, and NFA with $\varepsilon$-moves all share exactly the same expressive power. They all define the class of \textbf{Regular Languages}.

\section{Regular Expressions}

Instead of using automata, we can define languages using an algebraic formalism known as \textbf{Regular Expressions}. There are three primary classes of regular expressions to consider.

\subsection{Restricted Regular Expressions}

Restricted regular expressions are built from basic symbols (e.g., $a, b \in \mathcal{A}$) using exactly three operations:
\begin{itemize}
\item \textbf{Union ($\cup$)}: The set-theoretic union of two languages.
\item \textbf{Concatenation ($\cdot$)}: The concatenation of languages.
\item \textbf{Kleene Star ($^*$)}: The closure of a language (zero or more repetitions).
\end{itemize}

For example, $a \cdot b^*$ generates words like $a, ab, abb, abbb$, etc.

\subsection{General Regular Expressions}

General regular expressions extend restricted ones by adding two more set-theoretic operations:
\begin{itemize}
\item \textbf{Intersection ($\cap$)}: The intersection of two languages.
\item \textbf{Complementation ($\sim$)}: The complement with respect to the universal set $\mathcal{A}^*$.
\end{itemize}

\subsection{Star-free Regular Expressions}

Star-free regular expressions are general regular expressions but with the \textbf{Kleene Star ($^*$) operation forbidden}. They only permit union, concatenation, intersection, and complementation.

\subsection{Expressive Power of Regular Expressions}

The relationship between these three classes in terms of expressive power yields surprising results:

\begin{theorem}
Restricted regular expressions are exactly as expressive as general regular expressions. 
\end{theorem}
Adding intersection and complementation does not increase the class of languages we can define. Both restricted and general regular expressions define the exact same class of Regular Languages (equivalent to finite automata).

\begin{theorem}
Star-free regular expressions are strictly less expressive than general regular expressions.
\end{theorem}
Removing the Kleene star fundamentally reduces the expressive power, meaning there are regular languages that cannot be expressed without the Kleene star operation.
